\chapter*{Conclusion générale}\addcontentsline{toc}{chapter}{Conclusion générale}
Ce travail a relié deux problèmes souvent traités séparément : lire les caractères d'un document et comprendre assez de sa structure pour aider un analyste à préparer un dossier. La première tâche a fourni un contrat de preuves multi-pages où chaque élément garde sa source, sa page et ses coordonnées. La seconde a exploité ce contrat pour produire une hiérarchie de cahier des charges, des faits candidats, des bordereaux et des passages interrogables. L'espace de travail, le stockage local, la revue et le chiffrage séparé transforment ces résultats en parcours consultable.

La validation établit une base logicielle cohérente : 305 tests réussis lors de la vérification actuelle, des mesures de développement CDC explicitement bornées et un parcours local documenté. Elle établit aussi ce qui manque. Le seul BOQ réel de référence est vide ; les calculs vérifiés sur des fixtures synthétiques ne prouvent pas l'extraction de prix réels. Le test OCR avec modèles locaux n'a pas été rejoué pendant la rédaction, et la récupération de questions reste évaluée sur un petit ensemble contrôlé. Ces limites ne diminuent pas la valeur de la plateforme comme prototype d'ingénierie : elles définissent précisément le travail de données, de mesure et de contrôle nécessaire avant une utilisation engageante.

La contribution principale est la continuité de la preuve, du pixel ou du mot PDF jusqu'à la valeur présentée à l'utilisateur. Pour Aymen, le projet consolide des compétences concrètes d'ingénierie Data Science/IA : concevoir des contrats de données, intégrer OCR et extraction d'information, interpréter des annotations limitées, analyser les échecs et faire vérifier les sorties par un humain. Cette expérience rappelle qu'un système d'IA utile se juge aussi sur sa traçabilité et sa capacité à signaler ce qu'il ignore. Pour la suite, le progrès le plus décisif sera un corpus indépendant de CDC et de BOQ remplis, accompagné de protocoles de revue et de sécurité adaptés au contexte métier.
